\documentclass[12pt,a4paper]{amsart}

\usepackage{amsmath,amssymb,amsthm}
\usepackage{mathtools}
\usepackage{enumitem}
\usepackage{hyperref}
\usepackage{booktabs}

%%%─── Theorem environments
\newtheorem{theorem}{Theorem}[section]
\newtheorem{lemma}[theorem]{Lemma}
\newtheorem{proposition}[theorem]{Proposition}
\newtheorem{corollary}[theorem]{Corollary}
\newtheorem{conjecture}[theorem]{Conjecture}
\newtheorem{hypothesis}[theorem]{Hypothesis}
\theoremstyle{definition}
\newtheorem{definition}[theorem]{Definition}
\newtheorem{assumption}[theorem]{Assumption}
\newtheorem{problem}[theorem]{Open Problem}
\theoremstyle{remark}
\newtheorem{remark}[theorem]{Remark}
\newtheorem*{remark*}{Remark}

%%%─── Macros
\newcommand{\norm}[1]{\|#1\|}
\newcommand{\ip}[2]{\langle #1,\, #2\rangle}
\newcommand{\Kc}{\mathcal{K}_c}
\newcommand{\Kairy}{\mathcal{K}_{\mathrm{Ai}}}
\newcommand{\lambdac}[2]{\lambda_{#2}^{(#1)}}
\newcommand{\gapc}[2]{\mathrm{gap}_{#2}^{(#1)}}
\newcommand{\Ai}{\mathrm{Ai}}
\newcommand{\R}{\mathbb{R}}
\newcommand{\Det}{\mathrm{Det}}

\subjclass[2020]{47B34, 45C05, 35P20, 47A55, 60B20}
\keywords{Prolate spheroidal wave functions, transition regime, Airy kernel,
  Fredholm determinants, spectral universality, Widom asymptotics,
  sinc kernel, coalescing saddle point, Hilbert--P\'{o}lya programme}

\begin{document}

\title[Airy Universality for the PSWF Transition Regime]{%
  Airy Universality for the Prolate Concentration Operator in the
  Transition Regime: Fredholm Determinant Formulation and Spectral Corollaries}

\author{[Author]}
\address{[Address]}
\email{[Email]}

\begin{abstract}
This paper isolates the single asymptotic statement that underlies the sharp
transition-window spectral theory of the prolate concentration operator.
Instead of separating local model identification, quasimode construction,
and local counting into formally distinct assumptions, we formulate one
universality statement at the level of Fredholm determinants for the
appropriately rescaled transition-window operator.

\medskip
The central claim is that, in the regime $n\sim 2c/\pi$, the rescaled
Fredholm determinants of the prolate concentration operator converge to those
of the Airy kernel operator.
This places the PSWF transition regime in the same universality class as the
soft edge in random matrix theory.
Once this is available, the Airy-scale gap law, the local two-point counting
statement needed in Paper~XIII, and the existence of stable Airy quasimodes
all follow as corollaries.

\medskip
\textbf{Main output (skeleton).}
\begin{itemize}
\item \emph{(Theorem~U.)} Airy universality for the PSWF transition regime
  in Fredholm determinant form.
\item \emph{(Corollary~A.)} Airy-scale gap asymptotics $\gapc{c}{n}\asymp c^{-1/3}$.
\item \emph{(Corollary~B.)} Local eigenvalue counting
  $N_c([\mu_{n+1}(c)-\varepsilon(c),\mu_n(c)+\varepsilon(c)])=2$.
\item \emph{(Corollary~C.)} Existence of stable transition-window quasimodes.
\end{itemize}

\medskip
\textbf{Status.}
This paper is an open research skeleton.
Its purpose is to identify the exact analytic heart of the series:
a rigorous steepest-descent / integral-operator universality theorem for the
rescaled sinc kernel in the transition regime.
\end{abstract}

\maketitle

\section{Introduction}
\label{sec:intro}

Paper~XIII reduced the sharp transition-window gap problem to two visible
inputs: Airy-scale quasimodes and a local counting statement.
The present paper sharpens that reduction once more.
Its main point is that these two inputs are not analytically independent.
Both are consequences of a single deeper assertion: Airy universality for the
rescaled transition-window operator.

The correct formulation of this universality statement is not merely pointwise
kernel convergence.
Rather, the structurally stable object is the Fredholm determinant of the
rescaled operator.
Convergence at that level controls the local spectral statistics, the local
counting law, and the spectral projection structure needed to construct
quasimodes.
Thus the hard analytic core of the series is a single problem, not two.

\begin{center}
\fbox{\parbox{0.9\linewidth}{\centering\itshape
The transition regime of the PSWF concentration operator should lie in the
Airy universality class.
The correct rigorous statement is convergence of the Fredholm determinants of
the rescaled transition-window operator to those of the Airy kernel.
Gap asymptotics, counting, and quasimode stability are then consequences of
this one universality theorem.}}
\end{center}

\section{Operator setup and scaling}
\label{sec:setup}

\begin{definition}[Prolate concentration operator]
For $c>0$, let $\Kc:L^2([-1,1])\to L^2([-1,1])$ be the integral operator
\[
  (\Kc f)(x) := \int_{-1}^1 \frac{\sin(c(x-y))}{\pi(x-y)} f(y)\,dy.
\]
Its eigenvalues are denoted
\[
  1 > \lambdac{c}{0} \ge \lambdac{c}{1} \ge \lambdac{c}{2} \ge \cdots >0.
\]
\end{definition}

\begin{definition}[Transition regime]
Fix $A>0$.
The transition regime consists of indices
\begin{equation}\label{eq:transition-regime}
  n \in \Bigl[\frac{2c}{\pi}-Ac^{1/3},\,\frac{2c}{\pi}+Ac^{1/3}\Bigr].
\end{equation}
This is the natural window in which the prolate spectrum changes from bulk
behaviour to edge behaviour.
\end{definition}

\begin{definition}[Rescaled transition-window operator]
\label{def:rescaled-operator}
For $n$ in the transition regime, write $n = \frac{2c}{\pi}+\xi c^{1/3}$
with $\xi=O(1)$.
The rescaled transition-window operator $\widetilde{\Kc}^{(\xi)}$ is defined
by the kernel
\[
  \widetilde K_c^{(\xi)}(s,t)
  := c^{-1/3}\,K_c\bigl(\tfrac12 + c^{-1/3}s,\; \tfrac12 + c^{-1/3}t\bigr),
\]
where $s,t$ range over an interval $[-R,\infty)$ for a large cutoff $R>0$,
and where spatial and spectral scalings are coupled as in
Remark~\ref{rem:normalisation} below.
The exact normalisation is chosen so that the conjectural limiting kernel is
the Airy kernel on a half-line.
\end{definition}

\begin{remark}[Normalisation ansatz]
\label{rem:normalisation}
The centering is at the transition spectral level $\lambda=\tfrac12$.
With the index parametrisation $n = \frac{2c}{\pi}+\xi c^{1/3}$, the local
spectral variable is
\[
  s = c^{1/3}\Bigl(\lambdac{c}{n} - \tfrac12\Bigr),
\]
and the spatial rescaling near the boundary $x=1$ is $x=1-c^{-2/3}\tilde x$.
The rescaling coefficients are both of order $c^{1/3}$, i.e.\ $a_c,b_c\asymp c^{1/3}$,
chosen so that the cubic term in the local phase function survives in the
limit and produces the Airy kernel.
This is the normalisation underlying Theorem~\ref{thm:universality}.
\end{remark}

\begin{definition}[Airy kernel operator]
The Airy kernel is
\begin{equation}\label{eq:airy-kernel}
  K_{\mathrm{Ai}}(x,y)
  := \frac{\Ai(x)\Ai'(y)-\Ai'(x)\Ai(y)}{x-y},
\end{equation}
and the associated integral operator on $L^2([s,+\infty))$ is denoted by
$\Kairy^{(s)}$.
Its Fredholm determinant
\[
  F_{\mathrm{Ai}}(s;\lambda) := \Det(I-\lambda \Kairy^{(s)})
\]
encodes the soft-edge spectral statistics.
\end{definition}

\section{Universality theorem}
\label{sec:universality}

The entire paper is organised around the following single theorem.

\begin{theorem}[Airy universality in Fredholm determinant form]
\label{thm:universality}
Let $n$ lie in the transition regime \eqref{eq:transition-regime} and let
$\widetilde{\Kc}^{(\xi)}$ be the rescaled transition-window operator from
Definition~\ref{def:rescaled-operator}, with normalisation as in
Remark~\ref{rem:normalisation}.
Then for every compact set $\Lambda\subset\mathbb{C}$ and every compact
interval $S\subset\R$, one has
\begin{equation}\label{eq:fredholm-universality}
  \sup_{\lambda\in\Lambda}\sup_{s\in S}
  \bigl|\Det(I-\lambda\widetilde{\Kc}^{(\xi)}(s))
       - \Det(I-\lambda\Kairy^{(s)})\bigr|
  \longrightarrow 0
  \qquad (c\to\infty),
\end{equation}
uniformly for $\xi$ in compact sets.
\end{theorem}

\begin{remark}[Meaning of Theorem~\ref{thm:universality}]
The convergence \eqref{eq:fredholm-universality} is stronger than pointwise
kernel convergence.
It identifies the limiting determinantal process and therefore controls the
full local spectral statistics of the transition-window spectrum.
In particular, it encodes local counting distributions, consecutive spacing
statistics, and the stability of local spectral projections.
\end{remark}

\begin{lemma}[Trace-class criterion for determinant convergence]
\label{lem:trace-class}
To prove Theorem~\ref{thm:universality}, it is sufficient to establish that
for every compact interval $S\subset\R$,
\begin{equation}\label{eq:trace-class-conv}
  \sup_{s\in S}
  \norm{\widetilde{\Kc}^{(\xi)}(s)-\Kairy^{(s)}}_{\mathfrak{S}_1}
  \longrightarrow 0
  \qquad (c\to\infty),
\end{equation}
uniformly for $\xi$ in compact sets, where $\mathfrak{S}_1$ denotes
the trace-class norm.
\end{lemma}

\begin{proof}[Proof sketch]
Fredholm determinants are continuous with respect to the trace-class norm on
bounded subsets of trace-class operators.
Hence \eqref{eq:trace-class-conv} implies locally uniform convergence of the
corresponding determinants in the parameter $\lambda$.
This yields \eqref{eq:fredholm-universality}.
\end{proof}

\begin{remark}[Analytic status]
Theorem~\ref{thm:universality} is the main open theorem of the paper.
A proof would require a steepest-descent or equivalent asymptotic analysis of
the rescaled sinc kernel in the transition regime, together with a trace-class
control strong enough to pass from kernel asymptotics to Fredholm determinant
asymptotics.
Lemma~\ref{lem:trace-class} makes this bridge explicit.
This is the true analytic heart of the entire series.
\end{remark}

\section{Asymptotic kernel programme}
\label{sec:kernel-programme}

This section constitutes the first genuine analytic input:
we derive the phase structure of the sinc kernel in the transition regime,
identify the coalescing saddle point, and isolate a concrete error kernel
that can be estimated in Hilbert--Schmidt and trace norms.

\subsection{Integral representation and phase decomposition}

Using the Fourier representation of the sinc function,
\begin{equation}\label{eq:sinc-fourier}
  \frac{\sin(c(x-y))}{\pi(x-y)}
  = \frac{1}{2\pi}\int_{-c}^{c} e^{i\eta(x-y)}\,d\eta,
\end{equation}
the kernel of $\Kc$ becomes
\[
  K_c(x,y) = \frac{1}{2\pi}\int_{-c}^{c} e^{i\eta(x-y)}\,d\eta,
  \qquad x,y\in[-1,1].
\]
Introduce transition-window spatial variables centred near the boundary,
\[
  x = 1-c^{-2/3}s,
  \qquad
  y = 1-c^{-2/3}t,
  \qquad s,t\ge 0,
\]
and write the rescaled spectral parameter as
\[
  \xi = c^{1/3}\Bigl(\lambda-\frac12\Bigr).
\]
After the change of variables $\eta = c + c^{1/3}u$, the rescaled kernel
admits a local oscillatory representation of the form
\begin{equation}\label{eq:kernel-local-repr}
  \widetilde K_c^{(\xi)}(s,t)
  = \int_{\Gamma_c} \mathcal A_c(u;s,t,\xi)
    e^{i\Psi_c(u;s,t,\xi)}\,du
  \;+
  \mathcal E_c^{\mathrm{tail}}(s,t,\xi),
\end{equation}
where $\Gamma_c$ is the local contour, $\mathcal A_c$ is a smooth amplitude,
and $\mathcal E_c^{\mathrm{tail}}$ is the contribution of the complement of
the critical region.

\begin{definition}[Limiting phase and remainder]
\label{def:phase-remainder}
We decompose the phase as
\begin{equation}\label{eq:phase-split}
  \Psi_c(u;s,t,\xi) = \Psi_\infty(u;s,t,\xi) + R_c(u;s,t,\xi),
\end{equation}
with limiting cubic phase
\begin{equation}\label{eq:phase-infty}
  \Psi_\infty(u;s,t,\xi)
  := \frac{1}{3}u^3 + (s-t)u - \xi(s+t),
\end{equation}
and remainder term of the form
\begin{equation}\label{eq:phase-remainder-bound}
  R_c(u;s,t,\xi)
  = c^{-1/3}\Bigl(\alpha_1 u^2 + \alpha_2 u(s+t) + \alpha_3\Bigr)
    + c^{-2/3}Q_c(u;s,t,\xi),
\end{equation}
where $\alpha_1,\alpha_2,\alpha_3$ are bounded coefficients on compact
$(s,t,\xi)$-sets and $Q_c$ is polynomially bounded in $u,s,t$.
\end{definition}

\begin{remark}
The decomposition \eqref{eq:phase-split} is the correct starting point for
error analysis.
The target object is the two-variable Airy kernel $K_{\mathrm{Ai}}(s,t)$, not
the one-variable function $\Ai(s-t)$.
The latter appears only as an intermediate oscillatory profile; the full Airy
kernel emerges after the projection/half-line integration implicit in the
operator formulation.
\end{remark}

\subsection{Coalescing saddle point and cubic degeneracy}

The critical points of the limiting phase satisfy
\begin{equation}\label{eq:saddle-eq}
  \partial_u\Psi_\infty(u;s,t,\xi) = u^2 + (s-t) = 0.
\end{equation}
Hence the stationary points are
\[
  u_\pm = \pm i\sqrt{s-t},
\]
which coalesce at $u_*=0$ when $s=t$.
At the coalescence point one has
\[
  \partial_u^2\Psi_\infty(0;s,s,\xi)=0,
  \qquad
  \partial_u^3\Psi_\infty(0;s,s,\xi)=2\neq 0.
\]
This is exactly the Airy-type cubic degeneracy.

\begin{proposition}[Heuristic cubic phase reduction]
\label{prop:cubic-reduction}
On compact subsets of $(s,t,\xi)$-space, the local phase admits the expansion
\begin{equation}\label{eq:cubic-phase}
  \Psi_c(u;s,t,\xi)
  = \frac13 u^3 + (s-t)u - \xi(s+t)
    + c^{-1/3}\Bigl(\alpha_1 u^2 + \alpha_2 u(s+t) + \alpha_3\Bigr)
    + O\bigl(c^{-2/3}(1+|u|+|s|+|t|)^m\bigr)
\end{equation}
for some fixed exponent $m\ge 1$.
In particular, the quadratic term is lower order and the cubic term is the
leading non-trivial contribution at the coalescing saddle.
\end{proposition}

\begin{remark}[Analytic origin of the Airy kernel]
The Airy universality class is determined by the cubic degeneracy in
\eqref{eq:cubic-phase}.
The bandwidth edge $\eta=c$ of the sinc kernel plays the role of the soft
edge in random matrix theory, and the vanishing of the quadratic term at the
coalescence point is the phase-level mechanism behind the Airy limit.
\end{remark}

\subsection{Local error kernel}

We now isolate the error term that must be controlled.
Let
\begin{equation}\label{eq:error-kernel-def}
  \omega_c(s,t;\xi)
  := \bigl|\widetilde K_c^{(\xi)}(s,t) - K_{\mathrm{Ai}}(s,t)\bigr|.
\end{equation}
The analysis is based on a local/tail decomposition in the $u$-integral.
Fix $0<\delta<1/3$ and split
\begin{equation}\label{eq:local-tail-split}
  \Gamma_c = \Gamma_c^{\mathrm{loc}} \cup \Gamma_c^{\mathrm{tail}},
  \qquad
  \Gamma_c^{\mathrm{loc}} := \{u:|u|\le c^\delta\}.
\end{equation}
Then
\begin{equation}\label{eq:error-split}
  \omega_c(s,t;\xi)
  \le \omega_c^{\mathrm{loc}}(s,t;\xi)
     + \omega_c^{\mathrm{tail}}(s,t;\xi).
\end{equation}

\begin{lemma}[Variation estimate for the local error]
\label{lem:variation}
Assume the amplitude satisfies
\[
  |\mathcal A_c(u;s,t,\xi)-\mathcal A_\infty(u;s,t,\xi)|
  \le Cc^{-1/3}(1+|u|+|s|+|t|)^m
\]
on $\Gamma_c^{\mathrm{loc}}$, and let the phase decomposition be as in
Definition~\ref{def:phase-remainder}.
Then on compact subsets of $(s,t,\xi)$ one has
\begin{equation}\label{eq:variation-est}
  \omega_c^{\mathrm{loc}}(s,t;\xi)
  \le C\int_{\Gamma_c^{\mathrm{loc}}}
    \bigl(|R_c(u;s,t,\xi)| + c^{-1/3}(1+|u|+|s|+|t|)^m\bigr)
    e^{-\kappa (s+t)}\,|du|
\end{equation}
for some constants $C,\kappa>0$.
\end{lemma}

\begin{proof}[Proof sketch]
Write the difference between the local rescaled kernel and the limiting Airy
representation as amplitude error plus phase error.
For the phase part, use
\[
  e^{i(\Psi_\infty+R_c)} - e^{i\Psi_\infty}
  = e^{i\Psi_\infty}(e^{iR_c}-1)
\]
and the elementary estimate
\[
  |e^{iR_c}-1|\le |R_c|.
\]
This yields a Duhamel-type bound by integrating the absolute value of the
phase remainder and the amplitude discrepancy against the limiting oscillatory
weight.
The factor $e^{-\kappa(s+t)}$ records the standard Airy decay after contour
deformation to a steepest-descent path.
\end{proof}

\begin{corollary}[Model local kernel bound]
\label{cor:model-bound}
Under the hypotheses of Lemma~\ref{lem:variation}, there exist constants
$C,c_0>0$ and an integer $k\ge 0$ such that
\begin{equation}\label{eq:model-bound}
  \omega_c^{\mathrm{loc}}(s,t;\xi)
  \le C c^{-1/3}(1+|s|+|t|)^k e^{-c_0(s+t)}
\end{equation}
uniformly for $\xi$ in compact sets and $s,t\ge 0$.
\end{corollary}

\begin{proof}[Proof sketch]
Insert the explicit form \eqref{eq:phase-remainder-bound} of $R_c$ into
\eqref{eq:variation-est}.
The resulting integrals involve moments of $u$ against the Airy-type
oscillatory weight on $|u|\le c^\delta$.
These moments are bounded by a polynomial in $s,t$ times the exponential
Airy decay, giving \eqref{eq:model-bound}.
\end{proof}

\begin{lemma}[Tail estimate]
\label{lem:tail}
There exist constants $C,c_1>0$ such that
\begin{equation}\label{eq:tail-bound}
  \omega_c^{\mathrm{tail}}(s,t;\xi)
  \le C e^{-c_1 c^{\delta}} e^{-c_1(s+t)}
\end{equation}
uniformly for $\xi$ in compact sets and $s,t\ge 0$.
\end{lemma}

\begin{proof}[Proof sketch]
After contour deformation, the real part of the phase is strictly negative on
$\Gamma_c^{\mathrm{tail}}$.
The cubic term dominates there and produces exponential decay in $|u|$.
Since $|u|\ge c^\delta$ on the tail, the contribution is exponentially small
in $c^\delta$, uniformly in $s,t$ up to the standard Airy decay factor.
\end{proof}

\begin{proposition}[Global error kernel bound]
\label{prop:global-error}
Assume Lemma~\ref{lem:variation} and Lemma~\ref{lem:tail}.
Then there exist constants $C,c_0>0$ and an integer $k\ge 0$ such that
\begin{equation}\label{eq:global-error-bound}
  \omega_c(s,t;\xi)
  \le C c^{-1/3}(1+|s|+|t|)^k e^{-c_0(s+t)}
\end{equation}
for all sufficiently large $c$, uniformly for $\xi$ in compact sets and
$s,t\ge 0$.
\end{proposition}

\begin{proof}
Combine the local/tail decomposition \eqref{eq:error-split} with
Corollary~\ref{cor:model-bound} and Lemma~\ref{lem:tail}.
For large $c$, the exponential tail term is absorbed into the
$c^{-1/3}$ contribution.
\end{proof}

\subsection{Hilbert--Schmidt and trace-class upgrade}

\begin{lemma}[Hilbert--Schmidt control]
\label{lem:hs}
Assume Proposition~\ref{prop:global-error}.
Then
\begin{equation}\label{eq:hs-bound}
  \int_0^{\infty}\!\int_0^{\infty}
  \omega_c(s,t;\xi)^2\,ds\,dt
  \le C c^{-2/3}
\end{equation}
uniformly for $\xi$ in compact sets.
In particular,
\[
  \norm{\widetilde{\Kc}^{(\xi)}-\Kairy}_{\mathfrak S_2}
  \le C c^{-1/3}.
\]
\end{lemma}

\begin{proof}
Square the bound \eqref{eq:global-error-bound} and integrate over
$[0,\infty)^2$.
The exponential factor makes the integral finite, while the polynomial factor
is harmless.
\end{proof}

\begin{lemma}[Trace-class upgrade]
\label{lem:trace-upgrade}
Assume Proposition~\ref{prop:global-error}.
Then for every compact interval $S\subset\R$,
\begin{equation}\label{eq:trace-upgrade-bound}
  \sup_{s\in S}
  \norm{\widetilde{\Kc}^{(\xi)}(s)-\Kairy^{(s)}}_{\mathfrak S_1}
  \le C_S c^{-1/3}
\end{equation}
uniformly for $\xi$ in compact sets.
\end{lemma}

\begin{proof}[Proof sketch]
Use the exponential decay in \eqref{eq:global-error-bound} to factor the
kernel difference through weighted $L^2$ spaces.
The Hilbert--Schmidt bound from Lemma~\ref{lem:hs}, together with the
standard inequality $\norm{AB}_{\mathfrak S_1}
\le \norm{A}_{\mathfrak S_2}\norm{B}_{\mathfrak S_2}$ for an appropriate
factorisation, yields the trace-class estimate.
\end{proof}

\begin{problem}[From local phase analysis to determinant convergence]
\label{prob:traceclass-upgrade}
Make Proposition~\ref{prop:cubic-reduction} rigorous and prove the sequence
of estimates
\[
  \eqref{eq:variation-est}
  \Rightarrow
  \eqref{eq:model-bound}
  \Rightarrow
  \eqref{eq:global-error-bound}
  \Rightarrow
  \eqref{eq:trace-upgrade-bound}
  \Rightarrow
  \eqref{eq:fredholm-universality}.
\]
This is the full analytic route from coalescing saddle-point asymptotics to
Fredholm determinant universality.
\end{problem}

\section{Spectral corollaries}
\label{sec:corollaries}

Once Theorem~\ref{thm:universality} is available, the remaining results are
structural consequences.

\subsection{Gap scale}

\begin{corollary}[Airy-scale gap asymptotics]
\label{cor:gap}
Assume Theorem~\ref{thm:universality}.
Then for all sufficiently large $c$ and all transition-regime indices $n$,
\begin{equation}\label{eq:gap-asymp}
  \gapc{c}{n} \asymp c^{-1/3}.
\end{equation}
More precisely, after rescaling by $c^{1/3}$, consecutive PSWF gaps converge
to the corresponding consecutive Airy gaps.
\end{corollary}

\begin{proof}[Proof strategy]
Fredholm determinant convergence identifies the local limiting point process
with the Airy point process.
Consecutive eigenvalue spacings in the transition window therefore converge,
after the $c^{1/3}$ rescaling, to the Airy spacing law.
Since the Airy consecutive gaps are of order one, the unscaled PSWF gaps are
of order $c^{-1/3}$.
\end{proof}

\subsection{Local counting}

\begin{corollary}[Transition-window local counting]
\label{cor:counting}
Assume Theorem~\ref{thm:universality}.
Let $\mu_n(c)>\mu_{n+1}(c)$ be consecutive local model locations in the
transition regime, and let $\varepsilon(c)=o(c^{-1/3})$.
Then for all sufficiently large $c$,
\begin{equation}\label{eq:counting-cor}
  N_c\bigl([\mu_{n+1}(c)-\varepsilon(c),\,\mu_n(c)+\varepsilon(c)]\bigr)=2.
\end{equation}
In particular, the isolation condition required in Paper~XIII holds.
\end{corollary}

\begin{proof}[Proof strategy]
Fredholm determinant convergence implies convergence of the finite-dimensional
distributions of the associated determinantal point processes in the
transition window.
Hence the local counting statistics converge to those of the Airy point
process.
A window spanning exactly one consecutive Airy gap, enlarged by an
$o(c^{-1/3})$ perturbation, contains exactly two local eigenvalues for large
$c$.
Transporting this statement back through the transition scaling yields
\eqref{eq:counting-cor}.
\end{proof}

\subsection{Quasimodes}

\begin{corollary}[Stable Airy quasimodes]
\label{cor:quasimodes}
Assume Theorem~\ref{thm:universality}.
Then there exists a transition-window quasimode family
$\{f_n^{(c)}\}$ satisfying the quasimode hypotheses of Paper~XIII with
\begin{equation}\label{eq:quasimode-scale}
  \delta_{\mathrm{model}}(c) \asymp c^{-1/3},
  \qquad
  \varepsilon(c)=o(c^{-1/3}).
\end{equation}
Moreover the associated consecutive Gram matrices converge to the identity.
\end{corollary}

\begin{proof}[Proof strategy]
Fredholm determinant convergence implies convergence of local spectral
projections in the transition window.
Hence Airy eigenfunctions can be transported to approximate PSWF spectral
states with residual smaller than the Airy spacing scale.
The Gram convergence follows from the non-degeneracy of the limiting Airy
spectral projections and the simplicity of the local Airy spectrum.
\end{proof}

\section{Consequence for Paper~XIII}
\label{sec:paper13b}

The role of the present paper in the series is now completely transparent.
It supplies a single theorem whose corollaries verify all remaining open
inputs of Paper~XIII.

\begin{corollary}[Completion of the Paper~XIII transition argument]
\label{cor:paper13b}
Assume Theorem~\ref{thm:universality}.
Then all open transition-window inputs in Paper~XIII are verified.
Consequently, the quasimode transfer theorem of Paper~XIII applies and yields
\[
  \gapc{c}{n} \ge Cc^{-1/3}
\]
for all sufficiently large $c$ and all transition-regime indices $n$.
\end{corollary}

\begin{proof}
Corollary~\ref{cor:quasimodes} gives the quasimode package required in
Paper~XIII, and Corollary~\ref{cor:counting} gives the isolation/counting
input.
Applying Theorem~B and Corollary~D of Paper~XIII yields the claimed gap
bound.
\end{proof}

\section{Open problem and proof route}
\label{sec:open}

The previous sections show that the entire sharp transition-window theory has
been reduced to one analytic universality problem, whose proof programme is
now fully laid out in Section~\ref{sec:kernel-programme}.

\begin{problem}[Fredholm determinant universality for PSWF]
\label{prob:fredholm}
Prove Theorem~\ref{thm:universality} by executing the following chain:
\begin{enumerate}
\item Rigorous local phase decomposition
  \eqref{eq:phase-split}--\eqref{eq:phase-remainder-bound}.
\item Variation estimate for the local error kernel
  (Lemma~\ref{lem:variation}).
\item Global error kernel bound
  (Proposition~\ref{prop:global-error}).
\item Hilbert--Schmidt and trace-class upgrade
  (Lemmas~\ref{lem:hs} and \ref{lem:trace-upgrade}).
\item Passage to Fredholm determinant convergence
  via Lemma~\ref{lem:trace-class}.
\end{enumerate}
\end{problem}

\begin{remark}[Conceptual conclusion]
The sharp PSWF transition problem is a single universality problem.
The cubic degeneracy of the sinc phase at the bandwidth edge is the analytic
origin of the Airy universality class, and the decisive missing object is the
error kernel bound \eqref{eq:global-error-bound}.
Once that bound is proved, the determinant convergence machinery is standard.
\end{remark}

\begin{thebibliography}{99}

\bibitem{Slepian1961}
D.~Slepian, H.~O.~Pollak,
\textit{Prolate spheroidal wave functions, Fourier analysis and uncertainty~I},
Bell System Tech. J. \textbf{40} (1961), 43--63.

\bibitem{Widom1964}
H.~Widom,
\textit{Asymptotic behavior of the eigenvalues of certain integral equations},
Trans. Amer. Math. Soc. \textbf{109} (1964), 278--295.

\bibitem{Deift1999}
P.~Deift,
\textit{Orthogonal Polynomials and Random Matrices: A Riemann--Hilbert Approach},
Courant Lecture Notes, AMS, 1999.

\bibitem{TraWid1994}
C.~A.~Tracy, H.~Widom,
\textit{Level-spacing distributions and the Airy kernel},
Comm. Math. Phys. \textbf{159} (1994), 151--174.

\bibitem{Osipov2013}
A.~Osipov, V.~Rokhlin, H.~Xiao,
\textit{Prolate Spheroidal Wave Functions of Order Zero},
Springer, 2013.

\bibitem{Kato1966}
T.~Kato,
\textit{Perturbation Theory for Linear Operators},
Springer, 1966.

\bibitem{PaperVIII}
[Author], \textit{Paper~VIII}, 2026.

\bibitem{PaperXII}
[Author], \textit{Paper~XII}, 2026.

\bibitem{PaperXIII}
[Author], \textit{Paper~XIII (paper13b)}, 2026.

\end{thebibliography}

\end{document}
